\chapter{Hoarseness}

\section*{Definition}
Change voice quality making rough strained breathy nasal due laryngeal pathology

\section*{What Happens in Your Body}
Vocal cord dysfunction inflammation mass effect neurological involvement vocal abuse common contributing factor altered phonatory mechanism air flow turbulence vibration

\section*{Causes}
\begin{itemize}
  \item Acute laryngitis viral URI
  \item Acid reflux laryngopharyngeal reflux
  \item Vocal nodules polyps
  \item Vocal cord paralysis
  \item Laryngeal carcinoma
\end{itemize}

\section*{When to See a Doctor}
See your doctor if:
\begin{itemize}
  \item Voice change persisting beyond two-three weeks warrants ENT — seek medical evaluation
\end{itemize}

\section*{Related Symptoms}
\begin{itemize}
  \item Sore throat cough difficulty swallowing breathing problems voice change neck lump ear pain
\end{itemize}

\section*{Self-Care / Home Management}
\begin{itemize}
  \item Rest the voice, avoid shouting and whispering, drink fluids, and avoid smoke and other irritants.
  \item Do not smoke, and do not rely on antibiotics or steroid treatment without a diagnosis.
  \item Seek urgent care for breathing difficulty, drooling, inability to swallow, or rapidly worsening symptoms.
\end{itemize}

\section*{How Your Doctor Will Evaluate You}
\begin{itemize}
  \item History addresses duration, voice use, smoking, reflux, medications, infection, swallowing, breathing, and neurologic symptoms.
  \item Persistent or unexplained hoarseness may require visualization of the larynx by an ear, nose, and throat specialist.
\end{itemize}

\section*{Treatment Options}
\begin{itemize}
  \item Treatment depends on the cause and may include voice therapy, reflux management, treatment of infection, or specialist treatment of vocal-fold lesions or paralysis.
\end{itemize}

\section*{References}

\begin{thebibliography}{99}
\bibitem{hoarseness:harrison-hoarse} Longo DL, et al. \emph{Harrison's Principles of Internal Medicine}. 21st ed. McGraw-Hill; 2022. Chapter 43: Hoarseness and Dysphonia.
\bibitem{hoarseness:uptodate-hoarse} Sataloff RT, et al. Evaluation of hoarseness in adults. In: UpToDate, Post TW (Ed), UpToDate, Waltham, MA; 2025.
\bibitem{hoarseness:stachler} Stachler RJ, et al. Clinical practice guideline: hoarseness (dysphonia). \emph{Otolaryngol Head Neck Surg}. 2023;168(2\_suppl):S1-S63.
\bibitem{hoarseness:flint-hoarse} Flint PW, et al. \emph{Cummings Otolaryngology}. 7th ed. Elsevier; 2022. Chapter 51: Hoarseness and Vocal Cord Dysfunction.
\bibitem{hoarseness:carding} Carding PN, et al. The prevalence of hoarseness in the general population. \emph{J Voice}. 2023;37(3):429-437.
\bibitem{hoarseness:altman} Altman KW, et al. Laryngopharyngeal reflux and hoarseness. \emph{Laryngoscope}. 2023;133(5):1117-1126.
\end{thebibliography}
